\documentclass[reqno]{amsart}
\usepackage[english]{babel}
\usepackage{amscd,amssymb,amsmath,amsfonts,latexsym,amsthm}
\usepackage{inputenc}
\usepackage{graphicx}
\usepackage[fontsize=11pt]{scrextend}
\usepackage{hyperref,amsthm,amsmath,amssymb,color,multirow,graphicx}
\tolerance=5000 \topmargin -1cm \oddsidemargin=0,5cm
\evensidemargin=-0,2cm \textwidth 17cm \textheight 24cm
\linespread{1.0}
% ----------------------------------------------------------------
\vfuzz2pt % Don't report over-full v-boxes if over-edge is small
\hfuzz2pt % Don't report over-full h-boxes if over-edge is small

\newtheorem{thm}{Theorem}[section]
\newtheorem{cor}[thm]{Corollary}
\newtheorem{lem}[thm]{Lemma}
\newtheorem{prop}[thm]{Proposition}
\theoremstyle{definition}
\newtheorem{defn}[thm]{Definition}
\newtheorem{rem}[thm]{Remark}
\newtheorem{ex}[thm]{Example}
\renewcommand*{\proofname}{Proof:}

\numberwithin{equation}{section}


\def\abd{\mathop{\rm\leftharpoondown\!\!\rightharpoondown}}
\newcommand{\T}{\mathbb T}
\newcommand{\R}{\mathbb R}
\def\d{\,{\rm d}}
\def\E{\,{\rm e}}
\def\sgn{{\rm sgn}}
\def\sh{{\rm sh}}
\def\ch{{\rm ch}}
\def\abd{\mathop {\rm \leftharpoondown \!\!\!\!\! \rightharpoondown} }
\def\e{{\bf 1}\!\!{\rm I}}
\newcommand{\bbC}{\mathbb C}
\newcommand{\Ll}{L_2(\T^3)}
\newcommand{\cH}{\mathcal{H}}
\newcommand{\be}{\begin{equation}}

\newcommand{\ee}{\end{equation}}

%\usepackage{amssymb}

\renewcommand{\thefootnote}{\alph{footnote}}

\renewcommand{\thesection}{\arabic{section}}

\begin{document}


\sloppy

\begin{flushright}
	\begin{tabular}{ll}
		\textsf{Uzbek Mathematical Journal} &  \\
		\textsf{20??,\ Volume ??, Issue ?, pp.\pageref{firstpage}-\pageref{lastpage}}\\
		{DOI: 10.29229/uzmj.20??-?-??}\\
\end{tabular}
\end{flushright}



\begin{center}
\textbf{\large 7}\\

\textbf{ Mamanazarov Azizbek, Kodiraliev Asadbek}
\end{center}

\textbf{Abstract.}
In this paper, 

\textbf{Keywords:} 

\textbf{MSC (2020):35K70, 35D30, 35R11, 35A01}


\makeatletter
\renewcommand{\@evenhead}{\vbox{\thepage \hfil {\it Mamanazarov A., Kodiraliyev A.}  \hrule }}
\renewcommand{\@oddhead}{\vbox{\hfill
{\it Inverse coefficient problem }\hfill
\thepage \hrule}} \makeatother

\label{firstpage}



\section{Introduction}





Degenerate parabolic equations occupy a central place in the modern theory of partial
differential equations, both from a purely mathematical perspective and in view of their
numerous applications in mathematical physics and engineering. A particularly important
class of such equations involves diffusion operators that degenerate at the boundary of
the spatial domain, where the vanishing of the leading coefficient fundamentally alters
the analytical structure of the problem and demands the development of specialized
mathematical tools. Equations of this type arise naturally in the mathematical modeling of heat conduction and mass transfer in media with cylindrical or spherical symmetry \cite{Carslaw}. The study of boundary
value problems for such equations is of great significance not only from a theoretical
standpoint but also from the perspective of applications.
 
The spectral method based on Bessel function expansions, developed systematically for
equations with singular coefficients by Sabitov and his collaborators \cite{Sabitov,Vagapova},
constitutes an important analytical foundation for the present work. This approach
provides explicit solution representations as Fourier--Bessel series and has proved
effective for establishing existence, uniqueness and regularity of solutions to direct
boundary value problems for degenerate equations of the considered type. Direct and
inverse source problems for degenerate parabolic equations have been studied by a number
of authors. Hussein, Lesnic, Kamynin and Kostin \cite{A1} investigated both theoretical
and numerical aspects of such problems for equations with vanishing leading coefficients.
Cannarsa, Tort and Yamamoto \cite{A2} developed Carleman estimate techniques for
determining source terms in degenerate parabolic equations. Inverse problems for
time-dependent coefficients under integral overdetermination conditions were studied
in \cite{A3,A4}, where existence and uniqueness results were established via Volterra
integral equation arguments.
 
Alongside the classical integer-order theory, the analysis of time-fractional diffusion
equations has undergone remarkable development over the past two decades. The replacement
of the classical time derivative by the Caputo fractional derivative of order
$\alpha\in(0,1)$ gives rise to equations that model anomalous subdiffusion processes
exhibiting memory and hereditary effects, which cannot be captured by integer-order
models \cite{Podlubny}. The well-posedness theory for initial-boundary value problems for
time-fractional diffusion equations was developed by Luchko \cite{A5}, who established
existence, uniqueness and a maximum principle via eigenfunction expansion methods.
Boundary value problems for equations combining the Caputo fractional derivative and
the Bessel operator were studied by Agarwal, Karimov, Mamchuev and Ruzhansky \cite{A6},
where explicit solution representations were obtained by the spectral expansion method.
The study of inverse coefficient problems for time-fractional equations has also attracted
considerable attention; we refer to Durdiev \cite{A7} and Sakamoto--Yamamoto \cite{A8}
for results on the recovery of time-dependent coefficients from overdetermination
conditions.

 Inverse coefficient problems for fractional diffusion equations with the Bessel operator were studied by Akramova \cite{Akramova}, who considered the recovery of the lower-order coefficient in a one-dimensional fractional diffusion equation with the Bessel operator under an integral overdetermination condition, establishing existence via the Schauder principle and uniqueness via a Volterra integral equation argument.
 
A central problem in the theory of parabolic equations is the recovery of unknown
coefficients from indirect measurements of the solution  the so-called inverse
coefficient problem. Such problems are generally ill-posed in the sense of Hadamard,
and their rigorous treatment requires the proof of both existence and uniqueness of the
solution pair. A natural and effective approach is to supplement the governing equation
with an integral overdetermination condition, which permits the unknown coefficient to
be characterized as the solution of a Volterra integral equation of the second kind or
as a fixed point of a suitable operator \cite{Pedas,Zeidler}. Despite the growing body of
literature on inverse problems for both classical and time-fractional parabolic
equations, the case of degenerate equations with the Euler-Bessel operator with a general parameter $\nu$, particularly the simultaneous treatment of both classical and time-fractional settings within a unified framework, remains relatively unexplored.
 
In the present paper, we study direct and inverse coefficient problems in the domain
$\Omega = (0,l)\times(0,T)$ for two classes of degenerate parabolic equations:
\begin{equation}\label{1}
u_t = x^{-\nu}\frac{\partial}{\partial x}\left(x^{\nu}u_x\right) +r(t) f(x,t),\quad(x, t) \in \Omega=(0,l) \times(0, T),
\end{equation}
\begin{equation}\label{1.5}
    _c\mathcal{D}_t^\alpha u(x,t)=x^{-\nu}\frac{\partial}{\partial x}\left[x^{\nu}u_x\right]+q(t)u(x,t)+f(x,t), \quad (x, t) \in \Omega=(0,l) \times(0, T),
\end{equation}
where $r(t)$ and $q(t)$ are unknown time-dependent coefficients to be recovered, and
${}_{c}\mathcal{D}_t^{\alpha}$ denotes the Caputo fractional derivative of order
$\alpha\in(0,1)$ \cite{Podlubny}. For each equation, we prove the well-posedness of the
associated direct initial-boundary value problem and establish the existence and
uniqueness of the solution to the corresponding inverse problem, in which the unknown
coefficient is determined from an integral overdetermination condition. For
equation~\eqref{1}, the inverse problem is reduced to a Volterra integral equation
of the second kind, whose unique solvability follows from classical theory \cite{Pedas}.
For equation~\eqref{1.5}, existence is established via Schauder's fixed-point
theorem \cite{Zeidler}, and uniqueness is proved through an energy argument based on the
fractional Gronwall inequality \cite{Ye2007}.
 
The paper is organized as follows. Section~2 collects the necessary preliminary
material, including properties of the generalized Mittag-Leffler function, classical
functional inequalities, and the spectral problem associated with the degenerate
operator. Section~3 is devoted to the rigorous study of the direct and inverse problems
for equation~\eqref{1}. Section~4 addresses the well-posedness of the direct
problem for equation~\eqref{1.5}, and Section~5 is concerned with the corresponding
inverse coefficient problem.















\section{Preliminaries}

In this section, we recall some definitions and basic concepts that will be used throughout the paper. 
We will recall the definition of the generalized Mittag-Leffler function and give its necessary properties which we will use in our investigation. The
generalized Mittag-Leffler function is

\begin{equation*}
    E_{\alpha, \beta}(z)=\sum_{k=0}^{\infty} \frac{z^k}{\Gamma( \alpha k+\beta)}, z \in \mathbb{C} .
\end{equation*}

\begin{lem}\label{Mittag} \cite{Podlubny},  Suppose that $\alpha \in(0,2), \beta$ is an arbitrary real number and $\pi \alpha / 2<\mu<\min \{\pi, \pi \alpha\}$. Then there exists a positive constant $C$ such that the inequality holds:

\begin{equation*}
    \left|E_{\alpha, \beta}(z)\right| \leq \frac{C}{1+|z|},
\end{equation*}

for all $\mu \leq|\arg (z)| \leq \pi$ and $|z| \geq 0$.
\end{lem}
\begin{lem} 
 (Cauchy-Schwarz Inequality \cite{Evans}, p. 624]) Let $f$ and $g$ be measurable functions in $(0,l)$ such that $f, g \in L^2(0,l)$. Then
\begin{equation*}
    \left|\int_{0}^l f(x) g(x) \mathrm{d} x\right| \leq\left(\int_{0}^l|f(x)|^2 \mathrm{~d} x\right)^{1 / 2}\left(\int_{0}^l|g(x)|^2 \mathrm{~d} x\right)^{1 / 2}
\end{equation*}
\end{lem}
\begin{lem}\label{poincare}
Let $v\in C^1(0,l)$ and $v(0)=0$. Then there exists a constant 
$C_p=C_p(l,\nu)>0$ such that
\begin{equation}
\int_{0}^l x^{\nu}|v|^2\,dx
\le C_p
\int_{0}^l x^{\nu}|\partial_x v|^2\,dx .
\end{equation}
\end{lem}
\begin{proof}
    Let $\nu < 1$ be given. We have 
\begin{equation*}
    |v(x)|\leq \int_0^x y^{\nu/2}|\partial_y v|\frac{dy}{y^{\nu/2}}\leq \left(\int_0^xy^{-\nu}dy\right)^{1/2}\left(\int_0^x y^{\nu}|\partial_{y}v |^2dy\right)^{1/2}
\end{equation*}

\begin{equation*}
    =\left(\frac{x^{1-\nu}}{1-\nu}\right)^{1/2}\left(\int_0^x y^{\nu}|\partial_{y}v|^2dy\right)^{1/2}
\end{equation*}

Therefore, we obtain 

\begin{equation*}
    \int_0^l x^\nu |v|^2dx\leq \int_0^l \frac{x}{1-\nu}\int_0^x y^{\nu}|\partial_y v|^2dy=\frac{1}{1-\nu}\int_0^ly^\nu |\partial_y v|^2dy\int_y^lxdx
\end{equation*}

\begin{equation*}
    =\frac{1}{1-\nu}\int_0^l\frac{l^2-y^2}{2}y^{\nu}|\partial_y v|^2dy\leq\frac{l^2}{2(1-\nu)}\int_0^ly^\nu |\partial_y v|^2dy.
\end{equation*}
Then, we have 

\begin{equation*}
    \int_0^l x^\nu |v|^2dx\leq\frac{l^2}{2(1-\nu)}\int_0^lx^\nu |\partial_x v |^2dx.
\end{equation*}














    
\end{proof}
\begin{thm}\label{Schauder} 
(Schauder's Fixed-Point Theorem, \cite{Zeidler}) Let $X$ be a Banach space and $K \subset X$ a nonempty, closed, bounded, and convex subset. Suppose that the operator $A: K \rightarrow K$ is continuous and compact. Then the operator $A$ has at least one fixed point, i.e., there exists $x \in K$ such that

$$
A(x)=x.
$$
\end{thm} 

\begin{thm}\label{Gronwall}\cite{Ye2007} Let $\alpha\in(0,1)$, $b\geq 0$, and let 
$u(t),\,a(t):[0,T]\to[0,+\infty)$ be locally integrable functions. 
If
\begin{equation*}
    u(t)\leq a(t)+b\int_0^t(t-\tau)^{\alpha-1}u(\tau)\,d\tau,
    \qquad t\in[0,T],
\end{equation*}
then
\begin{equation*}
    u(t)\leq a(t)\cdot 
    E_{\alpha}\!\left(b\,\Gamma(\alpha)\cdot t^{\alpha}\right),
    \qquad t\in[0,T].
\end{equation*}
\end{thm}
\begin{rem}
 An operator $A: K \rightarrow K$ is called compact if it maps bounded subsets of $K$ to relatively compact subsets of $K$.
\end{rem} 
\begin{thm}\label{Ascoli}
 (Arzelà-Ascoli). Let $\mathcal{F} \subset C(\bar{\Omega})$, where $\bar{\Omega}=[0,l] \times[0, T]$. If $\mathcal{F}$ is uniformly bounded and equicontinuous in $\bar{\Omega}$, then every sequence in $\mathcal{F}$ has a uniformly convergent subsequence in $C(\bar{\Omega})$.
\end{thm}

\subsection{Spectral Problem}\label{sec5}

By applying Fourier's method to the Problem 1, we will have the following spectral problem: find the values of the parameter $\lambda$ for which there exist nontrivial solutions to the equation \eqref{1}.
For this aim, first, let  us consider the  following spectral problem 

\begin{equation}\label{5}
    x^{-\nu}\frac{\partial}{\partial x}\left[x^{\nu}X'(x)\right]=-\lambda^2 X(x),\quad x\in (0,l),
\end{equation}
satisfying the following conditions:
\begin{equation}\label{6}
X(0)=0, \qquad X(l)=0.
\end{equation}

We will show the existence of the eigenvalues and eigenfunctions of the problem $\left\{\eqref{5},\eqref{6}\right\}$. To do this,  we will first determine the general solution of \eqref{5}. By simplifying equation \eqref{5}, we obtain the following equation

\begin{equation}\label{eq5}
    X''(x)+\frac{\nu}{x}X'(x)+\lambda^2{X(x)}=0, \qquad x\in(0,l).
\end{equation}

In addition, with the substitution $X(x)=x^{(1-\nu)/2}Z(\xi)$ , $ \xi=\lambda x$, the last equation becomes the Bessel equation 

\begin{equation*}
    \xi^{2}Z''(\xi)+\xi{Z'(\xi)}+\left[\xi^2-p^2\right]Z(\xi)=0
\end{equation*}

where $p=\frac{1-\nu}{2}$.

Using the formula for the general solution of the Bessel equation \cite{bib6}, we will find the function $X(x)$ in the following  form 

\begin{equation}\label{7}
    X(x)=C_{1}x^{(1-\nu)/2}J_p(\lambda x)+C_2x^{(1-\nu)/2}Y_p(\lambda x)
\end{equation}

where $C_1$ and $C_2$ are arbitrary constants, $J_p(z)$ is the Bessel function of  the first kind of order $p$ defined by \cite{Watson}

\begin{equation*}
    J_p(z)=\sum_{n=0}^{+\infty} \frac{(-1)^n(z/2)^{2n+p}}{\Gamma(p+n+1)n!},
\end{equation*}

where $\Gamma(z)$ is the Euler's gamma function. 

In order for the function \eqref{7} to satisfy the first condition in equation \eqref{6}, we set  $C_2=0$ and, for convenience, take $C_1=1$, as eigenfunctions are determined only up to a multiplicative constant. we get the following form:

\begin{equation}\label{8}
    X(x)=x^{(1-\nu)/2}J_p(\lambda x).
\end{equation}

Then, obeying \eqref{7} to the condition \eqref{6}, we have the following equation 

\begin{equation}\label{9}
    J_p(\mu)=0, \qquad \mu=l\lambda.
\end{equation}

According to Lommel’s theorem \cite{Watson}, equation \eqref{9}, for $p>-1$, has a countable set of simple real roots 
$\mu_1<\mu_2<\mu_3<\dots<\mu_n<\cdots$ that determine the eigenvalues of the spectral problems \eqref{5} and \eqref{6}.
Setting $\mu=\mu_n=\lambda_nl$ in \eqref{8}, we obtain the corresponding eigenfunctions of the problem \eqref{5} and \eqref{6}:

\begin{equation*}
    \widetilde{X}_n(x)=x^{(1-\nu)/2}J_p(\lambda_nx), \quad \lambda_n=\frac{\mu_n}{l}, \quad n=1,2,...,
\end{equation*}
 
where $\mu_n$ is the $n$-th root of equation \eqref{8}.

For convenience in further computations, we orthonormalize this system of functions:

\begin{equation}\label{10}
    X_n(x)=\frac{1}{\left\|\widetilde{X}_n\right\|_{L_{2, \rho(0,l)}}}\widetilde{X}_n(x),
\end{equation}

where 

\begin{equation}\label{11}
    \left\|\widetilde{X}_n\right\|_{L_{2, \rho}(0,l)}^2=\int_0^l\rho(x) \widetilde{X}_n^2(x) d x, \quad \rho(x)=x^\nu.
    \end{equation}

It is known \cite{Watson} that the system of eigenfunctions \eqref{10} is complete in the space $L_2([0,l])$ with weight $x^\nu$.

 

Since the eigenvalues $\mu_n$ are positive zeros of the Bessel function
$J_p$, where $p=\frac{1-\nu}{2}$, we use the well-known asymptotic formula \cite{Olver}:
\begin{equation*}
    \mu_n=\pi\left(n+\frac{p}{2}-\frac14\right)+O\!\left(\frac{1}{n}\right),
\qquad n=1,2,3\dots,
\end{equation*}
which yields

\begin{equation}\label{12}
    \mu_n=\lambda_n l=\pi n-\frac{\nu}{4} \pi+O\left(\frac{1}{n}\right). \qquad n=1,2,3\dots.
\end{equation}


\section{The study of the problem 1}
In this section, we study the existence and uniqueness of solutions to the direct and inverse problems associated with equation \eqref{1} in the domain $\Omega$.

\subsection{Formulation of the direct problem 1}
In this section, in the domain $\Omega$, we will study the existence and uniqueness of the solution
of the following direct problem:\\
\textbf{Direct problem 1:} Throughout this section, we set $r(t)=1$ in Eq.~\eqref{1}, 
so that it reduces to
\begin{equation}\label{rt=1}
    u_t = x^{-\nu}\frac{\partial}{\partial x}\left(x^{\nu}u_x\right) + f(x,t),
    \quad (x,t)\in\Omega.
\end{equation}
Find a function $u(x, t)$ with the following regularity properties and conditions:

\begin{enumerate}
    \item $u(x,t), x^{\nu}u_x(x,t)\in C(\overline{\Omega})$, $x^{-\nu}\dfrac{\partial}{\partial x}\left(x^{\nu}u_x\right), u_t \in C(\Omega);$
    \item It satisfies Eq \eqref{rt=1} in the domain $\Omega$;
    \item At the boundary of the domain $\Omega$, it satisfies the following initial and boundary conditions:
\begin{equation}\label{2}
    u(x, 0)=\varphi(x), \quad x \in [0,l],
\end{equation}
\begin{equation}\label{3}
    u(0, t)=u(l, t)=0, \quad t \in[0, T] .
\end{equation}
$\varphi(x)$ and $q(t)$ are known functions, $\varphi(0)=\varphi(l)=0$  and $u(x,t)$ is an unknown function.
\end{enumerate}


\subsection{Existence and uniqueness of a solution to Direct Problem 1.}
We will prove the  unique solvability of Problem 1.
Let $u(x, t)$ be the solution of \eqref{rt=1}-\eqref{3}, 


Following \cite{Sabitov} and \cite{Vagapova}, we consider the functions

\begin{equation}\label{2.13}
    T_n(t)=\int_0^l u(x, t) x^\nu X_n(x) d x, \quad n=1,2, \ldots,
\end{equation}

where the $X_n(x)$ are defined by \eqref{10}. In view of \eqref{2.13} we introduce the functions

\begin{equation}\label{2.14}
    T_{n, \varepsilon}(t)=\int_{\varepsilon}^{l-\varepsilon} u(x, t) x^\nu X_n(x) d x, \quad n=1,2, \ldots,
\end{equation}


where $\varepsilon>0$ is sufficiently small. Differentiating \eqref{2.14} once with respect to $t \in (0,T)$ and using \eqref{rt=1}, we obtain that

$$T_{n, \varepsilon}^{\prime}(t)  =\int_{\varepsilon}^{l-\varepsilon} u_{t}(x, t) x^\nu X_n(x) d x$$

$$=\int_{\varepsilon}^{l-\varepsilon}\left(x^{-\nu} \frac{\partial}{\partial x}\left[x^\nu u_x\right]+f(x,t)\right) x^\nu X_n(x) d x $$

$$= \int_{\varepsilon}^{l-\varepsilon} \frac{\partial}{\partial x}\left(x^\nu u_x\right) X_n(x) d x +\int_{\varepsilon}^{l-\varepsilon}f(x,t)x^\nu X_n(x)dx$$




\begin{equation}\label{2.15}
  =\left.x^\nu u_x X_n(x)\right|_{\varepsilon} ^{l-\varepsilon}-\int_{\varepsilon}^{l-\varepsilon} x^\nu u_x X_n^{\prime}(x) d x +\int_{\varepsilon}^{l-\varepsilon}f(x,t)x^\nu X_n(x)dx.
\end{equation}


By \eqref{2.14}, in view of \eqref{eq5} we have
$$T_{n, \varepsilon}(t)  =-\frac{1}{\lambda_n^2} \int_{\varepsilon}^{l-\varepsilon} u(x, t) x^\nu\left[x^{-\nu}\frac{d}{d x} \left(x^\nu X_n^{\prime}(x)\right)\right] d x$$

$$=-\frac{1}{\lambda_n^2} \int_{\varepsilon}^{l-\varepsilon} u(x, t) \frac{d}{d x}\left(x^\nu X_n^{\prime}(x)\right) d x$$

\begin{equation}\label{2.16}
    = -\frac{1}{\lambda_n^2}\left[\left.u(x, t) x^\nu X_n^{\prime}(x)\right|_{\varepsilon} ^{l-\varepsilon}-\int_{\varepsilon}^{l-\varepsilon} x^\nu u_x X_n^{\prime}(x) d x\right].
\end{equation}


We find from \eqref{2.16} that

\begin{equation}\label{2.17}
    \int_{\varepsilon}^{l-\varepsilon} u_x x^\nu X_n^{\prime}(x) d x=\lambda_n^2 T_{n, \varepsilon}(t)+\left.u(x, t) x^\nu X_n^{\prime}(x)\right|_{\varepsilon} ^{l-\varepsilon}.
\end{equation}



Substituting \eqref{2.17} into \eqref{2.15}, we obtain

$$T_{n, \varepsilon}^{\prime }(t)=\left.x^\nu u_x X_n(x)\right|_{\varepsilon} ^{l-\varepsilon}-\lambda_n^2 T_{n, \varepsilon}(t)-\left.u(x, t) x^\nu X_n^{\prime}(x)\right|_{\varepsilon} ^{l-\varepsilon} $$

\begin{equation}\label{2.18}
    +\int_{\varepsilon}^{l-\varepsilon}f(x,t)x^\nu X_n(x)dx.
\end{equation}
Let $p=\frac{1-\nu}{2}$ and $z=\lambda_n x$. Using the identity
$\dfrac{d}{dz}\big(z^p J_p(z)\big)=z^p J_{p-1}(z)$ \cite{Watson}, we write
$$
X_n(x)=x^p J_p(\lambda_n x)=\lambda_n^{-p} z^p J_p(z).
$$
Hence,
$$
X_n'(x)=\lambda_n^{-p}\frac{d}{dx}\big(z^p J_p(z)\big)
=\lambda_n^{-p}\lambda_n z^p J_{p-1}(z)
=\lambda_n^{1-p}x^p J_{p-1}(\lambda_n x),
$$
that is
$$
X_n'(x)=\lambda_n^{(1+\nu)/2}x^{(1-\nu)/2}J_{-(1+\nu)/2}(\lambda_n x).
$$
Consequently,$\underset{x\to0+}{\lim}x^\nu X_n'(x) $ exists and is finite.

Then, letting $\varepsilon \rightarrow 0$ in \eqref{2.18} and using the boundary conditions \eqref{2}, \eqref{3}, and \eqref{6}, we obtain that $T_n(t)$ satisfies the differential equation
\begin{equation}\label{eq10}
T_n^{\prime}(t)+\lambda_{n}^{2} T_n(t)=f_n(t), \quad n=1,2,3, \ldots.
\end{equation}
The general solution of this equation is of the form:

\begin{equation}\label{2.20}
    T_n(t)=C_1 e^{-\lambda_n^2t}+\int_0^t f_n(\tau) e^{-\lambda_n^2\left(t-\tau\right)} d \tau, \quad n=1,2,3, \ldots,
\end{equation}
where $C_1$ is arbitrary constant.

We now use the boundary condition \eqref{2} and the formula \eqref{2.13} to find the
constant $C_1$:

\begin{equation}\label{2.23}
    T_n(0)=\int_0^l{u(x,0)x^\nu X_n(x)dx}=\int_0^l{\varphi(x)x^\nu X_n(x)dx=\varphi_n.}
\end{equation}
Substituting the function \eqref{2.20} into the boundary condition \eqref{2.23},
we obtain  $C_1=\varphi_n$. 

Using $C_1=\varphi_n$ and \eqref{2.20}, we finally find the function 


\begin{equation}\label{2.29}
     T_n(t)=\varphi_ne^{-\lambda_n^2t}+\int_0^t f_n(\tau) e^{-\lambda_n^2\left(t-\tau\right)} d \tau, \quad n=1,2,3, \ldots.
\end{equation}

Using \eqref{2.29} and \eqref{10}, we have the following expression for the function $u(x, t)$

\begin{equation}\label{15}
    u(x, t)=\sum_{n=1}^{+\infty}\left(\varphi_ne^{-\lambda_n^2t}+\int_0^t f_n(\tau) e^{-\lambda_n^2\left(t-\tau\right)} d \tau\right) X_n(x).
\end{equation}

\begin{lem}\label{lemm1}
The following bounds hold:
    \begin{equation}\label{eq1}
\left|T_n(t)\right|\le C_5\left(|\varphi_n|+\int_0^T|f_n(\tau)|
d\tau\right), \qquad 0\le \tau\le t\le T,    
    \end{equation}
    \begin{equation}\label{eq2}
        |T^{\prime}_n(t)|\le C_6\,n^2\big(|\varphi_n|+|f_n(t)|+\int_0^T|f_n(\tau)|d\tau\big), \qquad 0\le \tau\le t\le T.
    \end{equation}
Here and in what follows, $C_i$ are arbitrary constants.
\end{lem} 

\begin{proof}
Consider the function $T_n(t)$ defined by formula~\eqref{2.29}. Using \eqref{2.29} and the above estimate, we obtain
\begin{equation*}
    |T_n(t)|
\le |\varphi_n|\left|e^{-\lambda_n^2t}\right|
+\int_0^t|f_n(\tau)|
\left|e^{-\lambda_n^2\left(t-\tau\right)}\right|\,d\tau
\end{equation*}

$$
\le C_3|\varphi_n|+C_4\int_0^t |f_n(\tau)|\,d\tau
\le C_3|\varphi_n|+C_4\int_0^T |f_n(\tau)|\,d\tau .
$$
Consequently, estimate \eqref{eq1} follows:
$$
|T_n(t)|\le C_5\big(|\varphi_n|+\int_0^T|f_n(\tau)|d\tau\big),\qquad 0\le \tau\le t\le T.
$$

Next, from \eqref{eq10} we have
$$
T_n'(t)=f_n(t)-\lambda_n^2T_n(t).
$$
Therefore,
\begin{equation*}
    |T_n'(t)|
\le|f_n(t)|+\lambda_n^2\,|T_n(t)|.
\end{equation*}
Using \eqref{eq1} and the asymptotic bound $\lambda_n^2\le C n^2$ for large $n$, we obtain
$$
|T_n'(t)|\le C_6\,n^2\big(|\varphi_n|+|f_n(t)|+\int_0^T|f_n(\tau)|d\tau\big),\qquad t\in[0,T],
$$
which gives an estimate \eqref{eq2}.

The Lemma 1 has been proved.\end{proof}
\begin{lem}\label{lemm2}\cite{lemma3.2}     
For all $x \in[\delta, l]$, where $\delta$ is sufficiently small, we have
\begin{equation}\label{eq24}
    \left|X_n(x)\right| \leqslant C_7,
\end{equation}
\begin{equation}\label{eq25}
    \left|X_n^{\prime}(x)\right| \leqslant C_8 n,
\end{equation}

\begin{equation}\label{eq26}
     \left|X_n^{\prime \prime}(x)\right| \leqslant C_9 n^2.
\end{equation}

\end{lem}


\begin{lem}\label{lemm3}
    If $\varphi(x), f(x,t) \in C^4[0, l]$ and $\varphi(0)=f(0,t)=\varphi^{\prime}(0)=f^{\prime}(0,t)=\varphi^{\prime \prime}(0)= f^{\prime \prime}(0,t)=0, \varphi(l)=f(l,t)=\varphi^{\prime}(l)=f^{\prime}(l,t)=\varphi^{\prime \prime}(l)=f^{\prime \prime}(l,t)=0$, then we have
\begin{equation}\label{30}
    \left|\varphi_n\right| \leqslant \frac{C_{10}}{n^4}, \quad\left|f_n(t)\right| \leqslant \frac{C_{11}}{n^4} .
\end{equation}
\end{lem} 

 \begin{proof}
Integrating twice by parts in \eqref{2.23}, \eqref{eq10} and using \eqref{5}, \eqref{6} and the hypotheses of the lemma \ref{lemm3}, we have the following. 


$$\varphi_n  =\int_0^l \varphi(x) x^\nu X_n(x) d x=-\frac{1}{\lambda_n^2} \int_0^l \varphi(x)\frac{\partial}{\partial x}\left[x^{\nu}X_n^{\prime}(x)\right] d x $$

$$ =-\frac{1}{\lambda_n^2}\left[\left.\varphi(x) x^\nu X_n^{\prime}(x)\right|_0 ^l-\int_0^l \varphi^{\prime}(x) x^\nu X_n^{\prime} d x\right]=\frac{1}{\lambda_n^2} \int_0^l \varphi^{\prime}(x) x^\nu X_n^{\prime}(x) d x$$

$$  =\frac{1}{\lambda_n^2}\left[\left.\varphi^{\prime}(x) x^\nu X_n(x)\right|_0 ^l-\int_0^l\left(\varphi^{\prime}(x) x^\nu\right)^{\prime} X_n(x) d x\right]=-\frac{1}{\lambda_n^2} \int_0^l\left(\varphi^{\prime}(x) x^\nu\right)^{\prime} X_n(x) d x  $$

$$ =-\frac{1}{\lambda_n^2} \int_0^l \varphi^{\prime \prime}(x) x^\nu X_n(x) d x-\frac{\nu}{\lambda_n^2} \int_0^l \frac{\varphi^{\prime}(x)}{x} x^\nu X_n(x) d x.$$

We put in place


\begin{equation}\label{32}
    \varphi_n^{(2)}=\int_0^l \varphi^{\prime \prime}(x) x^\nu X_n(x) d x, \quad \varphi_{1 n}=\int_0^l \frac{\varphi^{\prime}(x)}{x} x^\nu X_n(x) d x.
\end{equation}

Then we obtain the representation

\begin{equation}\label{33}
    \varphi_n=-\frac{1}{\lambda_n^2} \varphi_n^{(2)}-\frac{\nu}{\lambda_n^2} \varphi_{1 n}.
\end{equation}



Under condition $\varphi^{\prime \prime}(0)=\varphi^{\prime \prime}(l)=0$ and in view of \eqref{5}, \eqref{6}, we integrate twice by parts in the first integral in \eqref{32}:


$$ \varphi_n^{(2)}=  \int_0^l \varphi^{\prime \prime}(x) x^\nu X_n(x) d x=-\frac{1}{\lambda_n^2} \int_0^l \varphi^{\prime \prime}(x) x^\nu\left[x^{-\nu}\frac{\partial}{\partial x}\left( x^{\nu}X_n^{\prime}(x)\right)\right] d x $$

$$=  -\frac{1}{\lambda_n^2} \int_0^l \varphi^{\prime \prime}(x)\left(x^\nu X_n^{\prime}(x)\right)^{\prime} d x=-\frac{1}{\lambda_n^2}\left[\left.\varphi^{\prime \prime}(x) x^\nu X_n^{\prime}(x)\right|_0 ^l\right. $$

$$ \left.-\int_0^l \varphi^{\prime \prime \prime}(x) x^\nu X_n^{\prime}(x) d x\right]=\frac{1}{\lambda_n^2} \int_0^l \varphi^{\prime \prime \prime}(x) x^\nu X_n^{\prime}(x) d x $$

$$ = \frac{1}{\lambda_n^2}\left[\left.\varphi^{\prime \prime \prime}(x) x^\nu X_n(x)\right|_0 ^l-\int_0^l\left(\varphi^{\prime \prime \prime}(x) x^\nu\right)^{\prime} X_n(x) d x\right]$$

$$ =  -\frac{1}{\lambda_n^2} \int_0^l \varphi^{(4)}(x) x^\nu X_n(x) d x-\frac{\nu}{\lambda_n^2} \int_0^l \frac{\varphi^{\prime \prime \prime}(x)}{x} x^\nu X_n(x) d x.$$



Introducing the notation

$$
\varphi_n^{(4)}=\int_0^l \varphi^{(4)}(x) x^\nu X_n(x) d x, \quad \varphi_{3 n}=\int_0^l \frac{\varphi^{\prime \prime \prime}(x)}{x} x^\nu X_n(x) d x,
$$

we obtain the representation

\begin{equation}\label{34}
    \varphi_n^{(2)}=-\frac{1}{\lambda_n^2} \varphi_n^{(4)}-\frac{\nu}{\lambda_n^2} \varphi_{3 n}.
\end{equation}

                   
Under condition $\varphi^{\prime}(0)=\varphi^{\prime}(l)=0$ and in view of \eqref{5}, \eqref{6}, we integrate twice by parts in the second integral in \eqref{32}:

$$ \varphi_{1 n}=  \int_0^l \frac{\varphi^{\prime}(x)}{x} x^\nu X_n(x) d x=-\frac{1}{\lambda_n^2} \int_0^l \frac{\varphi^{\prime}(x)}{x}\frac{\partial}{\partial x}\left(x^\nu X_n^{\prime}(x)\right) d x$$

$$ = -\frac{1}{\lambda_n^2} \int_0^l \varphi_1(x)\left[x^\nu X_n^{\prime}(x)\right]^{\prime} d x=-\frac{1}{\lambda_n^2}\left[\left.\varphi_1(x) x^\nu X_n^{\prime}(x)\right|_0 ^l\right.$$

$$ \left.-\int_0^l \varphi_1^{\prime}(x) x^\nu  X_n^{\prime}(x) d x\right]=\frac{1}{\lambda_n^2} \int_0^l \varphi_1^{\prime}(x) x^\nu X_n^{\prime}(x) d x$$

$$=  \frac{1}{\lambda_n^2}\left[\left.\varphi_1^{\prime}(x) x^\nu X_n(x)\right|_0 ^l-\int_0^l\left(\varphi_1^{\prime}(x) x^\nu \right)^{\prime} X_n(x) d x\right]$$

$$=  -\frac{1}{\lambda_n^2} \int_0^l\left(\varphi_1^{\prime}(x) x^\nu \right)^{\prime} X_n(x) d x=-\frac{1}{\lambda_n^2} \int_0^l \varphi_1^{\prime \prime}(x) x^\nu X_n(x) d x $$


\begin{equation}\label{35}
    -\frac{\nu}{\lambda_n^2} \int_0^l \frac{\varphi_1^{\prime}(x)}{x} x^\nu X_n(x) d x=-\frac{1}{\lambda_n^2} \varphi_{1 n}^{(2)}-\frac{\nu}{\lambda_n^2} \varphi_{2 n},
\end{equation}

where

$$
\varphi_{1 n}^{(2)}=\int_0^l \varphi_1^{\prime \prime}(x) x^\nu X_n(x) d x, \quad \varphi_{2 n}=\int_0^l \frac{\varphi_1^{\prime}(x)}{x} x^\nu X_n(x) d x, \quad \varphi_1(x)=\frac{\varphi^{\prime}(x)}{x}.
$$


Substituting \eqref{34} and \eqref{35} into \eqref{33}, we obtain the following

\begin{equation}\label{36}
    \varphi_n=\frac{1}{\lambda_n^4} \varphi_n^{(4)}+\frac{\nu}{\lambda_n^4} \varphi_{3 n}+\frac{\nu}{\lambda_n^4} \varphi_{1 n}^{(2)}+\frac{\nu^2}{\lambda_n^4} \varphi_{2 n} .
\end{equation}
 
The first bound in \eqref{30} follows from the representation \eqref{36}. The second bound in \eqref{30} is proved in a similar way. \end{proof} 

The solution of problems \eqref{rt=1}--\eqref{3} can be represented as a
Fourier--Bessel series in terms of the eigenfunctions \eqref{10} with
coefficients \eqref{2.29}:
\begin{equation}\label{37}
u(x,t)=\sum_{n=1}^{\infty}T_n(t)\,X_n(x).
\end{equation}

The formal termwise differentiation of \eqref{37} yields
\begin{equation}\label{39}
u_t(x,t)=\sum_{n=1}^{\infty}T_n'(t)\,X_n(x),\qquad
u_x(x,t)=\sum_{n=1}^{\infty}T_n(t)\,X_n'(x).
\end{equation}

Moreover, using the eigenvalue relation \eqref{5}, we obtain
\begin{equation}\label{opseries}
x^{-\nu}\frac{\partial}{\partial x}\bigl(x^\nu u_x(x,t)\bigr)
=\sum_{n=1}^{\infty}T_n(t)\,
x^{-\nu}\frac{d}{dx}\bigl(x^\nu X_n'(x)\bigr)
=-\sum_{n=1}^{\infty}\lambda_n^2 T_n(t)\,X_n(x).
\end{equation}

For every $(x,t)\in\bar{\Omega}$, the series
\eqref{37} is majorized by
\begin{equation}\label{maj1}
C_{12}\sum_{n=1}^{\infty}\bigl(|\varphi_n|+|f_n(t)|\bigr),
\end{equation}
while the series in \eqref{39} are majorized for $(x,t)\in\overline{\Omega}_\delta$
by
\begin{equation}\label{maj2}
\tilde{C}_{13}\sum_{n=1}^{\infty}\,n^2\big(|\varphi_n|+|f_n(t)|+\int_0^T|f_n(\tau)|d\tau\big), \qquad  C_{13}\sum_{n=1}^{\infty}
n\left(|\varphi_n|+\int_0^T|f_n(\tau)|
d\tau\right).
\end{equation}

Furthermore, the operator series \eqref{opseries} is majorized in
$\overline{\Omega}_\delta$ by
\begin{equation}\label{maj3}
C_{14}\sum_{n=1}^{\infty}
\lambda_n^2\left(|\varphi_n|+\int_0^T|f_n(\tau)|
d\tau\right),
\end{equation}
where $\Omega_\delta=\Omega\cap\{x>\delta\}$.

By Lemma~\ref{lemm3}, there exists a constant $C>0$ such that
$|\varphi_n|+|f_n(t)|\le C n^{-4}$.
Moreover, using the asymptotic relation \eqref{12}, we have
$\lambda_n^2\le C n^2$ for all sufficiently large $n$.
Therefore, the series in \eqref{maj1}--\eqref{maj3} are respectively majorized by
the convergent numerical series
\begin{equation*}
    \tilde{C}_{1}\sum_{n=1}^{\infty}n^{-4},\qquad
\tilde{C}_{2}\sum_{n=1}^{\infty}n^{-3},\qquad
\tilde{C}_{3}\sum_{n=1}^{\infty}n^{-2}.
\end{equation*}


Hence, the series \eqref{37}, \eqref{39}, and \eqref{opseries} converge uniformly
in their corresponding domains.


In the next step, we prove that the solution of problem \eqref{rt=1}--\eqref{3} in the form \eqref{15} is unique for any arbitrary function $r(t)\in C[0,l]$.

To this end, suppose that problem \eqref{rt=1}--\eqref{3} admits two solutions. Define
$
U(x,t)=u_1(x,t)-u_2(x,t).
$

Then $U(x,t)$ satisfies:
\begin{equation}\label{wtenglama}
U_t-x^{-\nu}\frac{\partial}{\partial x}\left[x^{\nu}U_x\right]=0
\end{equation}
\begin{equation}\label{wbshart}
    U(x,0)=0, \quad x\in[0,l]
\end{equation}
\begin{equation}\label{wchshart}
    U(0,t)=U(l,t)=0,\quad t\in[0,T]
\end{equation}

consider the energy functional: 

\begin{equation*}
    E(t)=\frac{1}{2}\int_0^l x^{\nu} |U|^2dx.
\end{equation*}
Take the derivative with respect to $t$ and then insert the PDE:
 \begin{equation*}
     \frac{d}{dt}E(t)=\int_0^lx^{\nu}UU_tdx=\int_0^lx^{\nu}U\left[x^{-\nu}\frac{\partial}{\partial x}\left(x^{\nu}U_x\right)\right]dx=\int_0^l U\frac{\partial}{\partial x} \left(x^{\nu}U_x\right)dx.
 \end{equation*}


Applying integration by parts and using the boundary conditions  $U(0,t)=U(l,t)=0$, we obtain:


\begin{equation*}
    \int_0^l U\frac{\partial}{\partial x} \left(x^{\nu}U_x\right)dx=\left.x^{\nu} U U_x\right|_0 ^l -\int_0^l x^{\nu} |U_x|^2dx=-\int_0^l x^{\nu} |U_x|^2dx,
\end{equation*}

then, 

\begin{equation*}
    \frac{d}{dt}E(t)=-\int_0^l x^{\nu} |U_x|^2dx.
\end{equation*}

Using the Lemma~\ref{poincare}, we get:

\begin{equation*}
    \int_0^l x^\nu|U|^2 d x \leq C_p \int_0^l x^\nu\left|U_x\right|^2 d x
\end{equation*}

By substituting the expression into the inequality, it follows that:

\begin{equation*}
    \frac{d}{dt}E(t)\leq-\frac{1}{C_p}\int_0^l x^{\nu} |U|^2dx.
\end{equation*}

Since $\int_0^lx^{\nu}|U|^2dx=2E(t)$, we have: 

\begin{equation*}
    \frac{d}{dt}E(t)\leq -\frac{2}{C_p}E(t).
\end{equation*}
Let $P=-\dfrac{2}{C_p}$. Then 

\begin{equation*}
    \frac{d}{dt}E(t)\leq PE(t).
\end{equation*}
Since $E(0)=0$, Gronwall's lemma gives:
\begin{equation*}
    E(t) \leq E(0)\, e^{Pt} = 0, \quad \text{for all } t\in[0,T].
\end{equation*}
On the other hand, by definition, $E(t) \geq 0$. 
Therefore, we conclude that $E(t) = 0$ for all $t\in[0,T]$.

Thus, $E(t)$ vanishes for all $t\in[0,T]$, which implies that $U(x,t)=0$ everywhere in $\Omega$. Therefore,
\begin{equation*}
    u_1(x,t)=u_2(x,t),
\end{equation*}
and the uniqueness follows.















 We have proved the following theorem:

\begin{thm}\label{theorem2}
Let $\varphi(x)$ and $f(x,t)$ satisfy the assumptions of Lemma~\ref{lemm3}.
Then there exists a unique solution $u(x, t)$ to Problem 1.
\end{thm}



\subsection{Existence and uniqueness of the solution of the inverse problem~1.}
In this section, we will show that the result of the inverse problem 1  is exist and unique.


Note that the solution representation \eqref{15}, obtained in the previous 
section for the case $r(t) \equiv 1$, can be naturally extended to the case 
of a general function $r(t) \in C[0,T]$ by replacing $1$ with $r(\tau)$ 
under the integral sign, yielding
\begin{equation}\label{15r}
    u(x,t) = \sum_{n=1}^{+\infty} \varphi_n e^{-\lambda_n^2 t} X_n(x)
    + \sum_{n=1}^{+\infty} \int_0^t r(\tau) f_n(\tau) 
    e^{-\lambda_n^2(t-\tau)}\, d\tau \, X_n(x).
\end{equation}
This representation will be used throughout this section.

\textbf{Problem 2.} Show the existence and uniqueness of a pair of functions $u(x, t), r(t)$ such that $r(t) \in C[0, T]$, and $u(x, t)$ is solution to the problem $\{\eqref{1}, \eqref{2},\eqref{3}\}$, satisfying the over-determination condition

\begin{equation}\label{1.6}
    \int_0^l u(x, t)x^{\nu} w(x) d x=g(t), \quad t \in[0, T].
\end{equation}
where $w(x)$ and $g(t)$ are given functions.
 Since the series \eqref{15r} and its first derivative with respect to $t$ converge uniformly in $\overline{\Omega}$, we may differentiate it term by term. Differentiating \eqref{15r} with respect to $t$, we obtain the following relation

\begin{equation*}
    u_t=-\sum_{n=1}^{+\infty} \lambda_n^2 \varphi_n e^{-\lambda_n^2 t} X_n(x)-\sum_{n=1}^{+\infty}\left(\int_0^t \lambda_n ^2r(\tau) f_n(\tau) e^{-\lambda_n^2(t-\tau)} d \tau-r(t) f_n(t)\right) X_n(x).
\end{equation*}

Hence, from the over-determination condition \eqref{1.6}, we have

$$
\begin{aligned}
\int_0^lx^{\nu} w(x) u_t(x, t) d x & =\sum_{n=1}^{+\infty}-\lambda_n^2 \varphi_n e^{-\lambda_n2 t} \int_0^lx^{\nu} w(x) X_n(x) d x+r(t) \sum_{n=1}^{+\infty} f_n(t) \int_0^l x^{\nu}w(x) X_n(x) d x- \\
& -\sum_{n=1}^{+\infty} \lambda_n^2 \int_0^t r(\tau) f_n(\tau) e^{-\lambda_n^2(t-\tau)} d \tau \int_0^l x^{\nu}w(x) X_n(x) d x=g^{\prime}(t)
\end{aligned}
$$

or

$$
\begin{gathered}
r(t)=\frac{g^{\prime}(t)+\sum_{n=1}^{+\infty} \lambda_n ^2\varphi_n e^{-\lambda_n^2 t} \int_0^l x^{\nu}w(x) X_n(x) d x}{\int_0^l w(x) f(x, t) d x}+ \\
+\frac{1}{\int_0^l w(x) f(x, t) d x} \int_0^t\left(\sum_{n=1}^{+\infty} \lambda_n^2 f_n(\tau) e^{-\lambda_n^2(t-\tau)} d \tau \int_0^l x^{\nu}w(x) X_n(x) d x\right) r(\tau) d \tau.
\end{gathered}
$$
This is equivalent to the following Volterra
integral equation of the second kind:
\begin{equation}\label{volterra_r}
r(t)=A(t)+\int_{0}^{t}K(t,\tau)r(\tau)\,d\tau,
\end{equation}

where

\begin{equation}\label{A}
    A(t) =
\frac{
g'(t)+\sum_{n=1}^{+\infty}\lambda_n^2\varphi_n e^{-\lambda_n^2 t}
\int_{0}^{l} x^{\nu}w(x)X_n(x)\,dx
}{
\int_{0}^{l} w(x)f(x,t)\,dx
},
\end{equation}
\begin{equation}\label{K}
   K(t,\tau) =
\frac{
\sum_{n=1}^{+\infty}\lambda_n^2 f_n(\tau)
e^{-\lambda_n^2 (t-\tau)}
\int_{0}^{l} x^{\nu}w(x)X_n(x)\,dx
}{
\int_{0}^{l} w(x)f(x,t)\,dx
}.
\end{equation}
Now, let us define the conditions under which the series in \eqref{A} and \eqref{K} converges.

\begin{equation*}
    \sum_{n=1}^{+\infty}\lambda_n^2\varphi_n e^{-\lambda_n^2 t}
\int_{0}^{l} w(x)X_n(x)\,dx\quad\text{and} \quad \sum_{n=1}^{+\infty}\lambda_n^2 f_n(\tau)
e^{-\lambda_n^2 (t-\tau)}
\int_{0}^{l} w(x)X_n(x)\,dx.
\end{equation*}





Taking into account the Lemma~\ref{lemm3}--\ref{lemm2}, $w(x) \in C^4[0, l]$, $\lambda_n^2\leq Cn^2$ and  Using the Cauchy-Schwarz inequality, we have 

\begin{equation*}
    \sum_{n=1}^{+\infty}\left|\lambda_n^2\varphi_n e^{-\lambda_n^2 t}
\int_{0}^{l} x^{\nu}w(x)X_n(x)\,dx\right| = C_{15} \sum_{n=1}^{+\infty}\lambda_n^2|\varphi_n||w_n|
\end{equation*}

\begin{equation*}
  \leq C_{15}\left(\sum_{n=1}^{+\infty}\lambda_n^2|\varphi_n|^2\right)^{1/2}\left(\sum_{n=1}^{+\infty}\lambda_n^2|w_n|^2\right)^{1/2}\leq \tilde{C}_{15}\left(\sum_{n=1}^{+\infty}n^{-2}\right)^{1/2}\left(\sum_{n=1}^{+\infty}n^{-2}\right)^{1/2}<+\infty
\end{equation*}
and 
\begin{equation*}
    \sum_{n=1}^{+\infty}\left|\lambda_n^2 f_n(\tau)
e^{-\lambda_n^2 (t-\tau)}
\int_{0}^{l} x^{\nu}w(x)X_n(x)\,dx\right|=C_{16} \sum_{n=1}^{+\infty}\lambda_n^2|f_n||w_n|
\end{equation*}

\begin{equation*}
    \leq C_{16}\left(\sum_{n=1}^{+\infty}\lambda_n^2\left|f_n(\tau)\right|^2\right)^{1/2}\left(\sum_{n=1}^{+\infty}\lambda_n^2\left|w_n\right|^2\right)^{1/2}\leq \tilde{C}_{16}\left(\sum_{n=1}^{+\infty}n^{-2}\right)^{1/2}\left(\sum_{n=1}^{+\infty}n^{-2}\right)^{1/2}<+\infty.
\end{equation*}



 Thus, the Volterra integral equation's last equation has a unique continuous solution.




\begin{thm}\label{3.5}
   Let the following conditions hold:
 $\varphi(x)$ and $f(x,t)$ satisfy the assumptions of Lemma~\ref{lemm3}; $w(x) \in C^{4}[0,l]$ with $w(0) = w(l) = 0$; $g \in C^{1}([0,T])$ with $g(t) \neq 0$ for all $t \in [0,T]$;  $\displaystyle\int_{0}^{l} w(x)f(x,t)\,dx \neq 0$ for all $t \in [0,T]$.
Then the inverse problem \{\eqref{1},\eqref{2},\eqref{3},\eqref{1.6}\} admits a unique solution
pair $\{r(t),\,u(x,t)\}$ such that
\begin{equation*}
    r(t) \in C([0,T]),\quad u(x,t),\quad x^{\nu}u_{x}(x,t) \in C(\overline{\Omega}),\quad x^{-\nu}\frac{\partial}{\partial x}\bigl[x^{\nu}u_{x}\bigr],\quad
    u_{t} \in C(\Omega).
\end{equation*}

\end{thm}



\section{The study of the problem 2}

In this section, in the domain $\Omega$, we will study the existence and uniqueness of the solution
of the following direct and inverse problem for Eq. \eqref{1.5}:

\subsection{ Formulation of the direct problem 2}

Find a function $ u(x, t)$ with the following properties:

\begin{enumerate}
    \item $u(x,t), x^{\nu}u(x,t)\in C(\overline{\Omega})$, $x^{-\nu}\dfrac{\partial}{\partial x}\left(x^{\nu}u_x\right), {_c}D_{0t}^{\alpha}u \in C(\Omega);$
    \item It satisfies Eq. \eqref{1.5} in the domain $\Omega$;
    \item On the boundary of the domain $\Omega$, it satisfies the initial \eqref{2} and the boundary conditions \eqref{3}:
\end{enumerate}
\subsection{Existence and Uniqueness Results of the Direct Problem 2}
In this subsection, we
present the results of the existence and uniqueness of the direct
problem \eqref{1.5},\eqref{2},\eqref{3}.
The solution is constructed through Fourier series. Consider the eigenfunctions $X_n(x)$ \eqref{10} and the eigenvalues $\lambda_n=\dfrac{\mu_n}{l},\quad \mu_n=\pi n-\dfrac{\nu}{4}\pi+O(\dfrac{1}{n})$.
Expand:

$$
u(x, t)=\sum_{n=1}^{\infty} T_n(t) X_n(x), \quad \varphi(x)=\sum_{n=1}^{\infty} \varphi_n X_n(x), \quad f(x, t)=\sum_{n=1}^{\infty} f_n(t) X_n(x),
$$

where $\varphi_n=\left\langle\varphi, X_n\right\rangle,\quad f_n(t)=\left\langle f(\cdot, t), X_n\right\rangle$. Substitution into the PDE yields:

\begin{equation}\label{T(t)eq}
    _cD_{0t}^\alpha T_n(t)+\lambda_n^2T_n(t)=q(t)T_n(t)+f_n(t), \quad T_n(0)=\varphi_n.
\end{equation}

Using the definition of the Caputo fractional derivative, equation \eqref{T(t)eq} is equivalent to the following Volterra integral equation:

\begin{equation}\label{4.1}
  T_n(t)
=
\varphi_n E_{\alpha}\!\left(-\lambda_n^2 t^{\alpha}\right)
+
\int_0^t
(t-\tau)^{\alpha-1}
E_{\alpha,\alpha}\!\left(-\lambda_n^2 (t-\tau)^{\alpha}\right)
\big[q(\tau)T_n(\tau)+f_n(\tau)\big]
\, d\tau.
\end{equation}

\begin{lem}\label{T(t)}
    According to the assumptions of Theorem~\ref{Gronwall}, \ref{Mittag},lemma~\ref{lemm3} and $q(t)\in C[0,T]$, the solution $T_n(t)$ of
\eqref{4.1} satisfies the estimate
\begin{equation*}
 |T_n(t)|\leq  \frac{C_{}}{n^4}.
\end{equation*}
\end{lem}
\begin{proof}
Consider the function $T_n(t)$ defined by the Volterra integral 
equation~\eqref{4.1}:
\begin{equation*}
T_n(t) = \varphi_n E_{\alpha}\!\left(-\lambda_n^2 t^{\alpha}\right)
+ \int_0^t (t-\tau)^{\alpha-1}
  E_{\alpha,\alpha}\!\left(-\lambda_n^2(t-\tau)^{\alpha}\right)
  \bigl[q(\tau)T_n(\tau)+f_n(\tau)\bigr]\,d\tau.
\end{equation*}


Since $\alpha\in(0,1)$ and $\lambda_n^2 t^\alpha \geq 0$, by 
Lemma~\ref{Mittag} there exists a constant $C_1>0$ such that
\begin{equation}\label{E<C}
    \left|E_{\alpha}\!\left(-\lambda_n^2 t^{\alpha}\right)\right| 
    \leq C_1, 
    \qquad
    \left|E_{\alpha,\alpha}\!\left(-\lambda_n^2(t-\tau)^{\alpha}\right)\right| 
    \leq \frac{C_1}{1+\lambda_n^2(t-\tau)^{\alpha}} \leq C_1,
\end{equation}
for all $0\leq\tau\leq t\leq T$.

Since $q\in C[0,T]$, there exists $M>0$ such that 
$|q(t)|\leq M$ for all $t\in[0,T]$. 
Taking absolute values in~\eqref{4.1} and applying \eqref{E<C}:
\begin{align*}
|T_n(t)| 
&\leq |\varphi_n|\left|E_{\alpha}\!\left(-\lambda_n^2 t^{\alpha}\right)\right|
 + \int_0^t (t-\tau)^{\alpha-1}
   \left|E_{\alpha,\alpha}\!\left(-\lambda_n^2(t-\tau)^{\alpha}\right)\right|
   |q(\tau)|\,|T_n(\tau)|\,d\tau \\
&\quad + \int_0^t (t-\tau)^{\alpha-1}
   \left|E_{\alpha,\alpha}\!\left(-\lambda_n^2(t-\tau)^{\alpha}\right)\right|
   |f_n(\tau)|\,d\tau \\[4pt]
&\leq C_1|\varphi_n|
 + C_1 M\int_0^t (t-\tau)^{\alpha-1}|T_n(\tau)|\,d\tau
 + C_1\int_0^t (t-\tau)^{\alpha-1}|f_n(\tau)|\,d\tau.
\end{align*}

Define:
\begin{equation}\label{a_t}
    a(t) = C_1|\varphi_n| 
           + C_1\int_0^t (t-\tau)^{\alpha-1}|f_n(\tau)|\,d\tau,
    \qquad b= C_1 M \geq 0.
\end{equation}
Then the inequality takes the form:
\begin{equation}\label{weakly_singular}
    |T_n(t)| \leq a(t) + b\int_0^t (t-\tau)^{\alpha-1}|T_n(\tau)|\,d\tau.
\end{equation}


Note that $a(t)$ defined in~\eqref{a_t} is non-negative and 
non-decreasing on $[0,T]$. 
Applying Lemma~\ref{Gronwall} to~\eqref{weakly_singular} 
with $\beta=\alpha$ and constant $b=C_1 M$, we obtain:
\begin{equation}\label{Tn_ML}
    |T_n(t)| \leq a(t)\cdot 
    E_{\alpha}\!\left(C_1 M\,\Gamma(\alpha)\cdot t^{\alpha}\right),
    \qquad t\in[0,T].
\end{equation}
Since $\alpha\in(0,1)$, $T<\infty$, and $M$, $C_1$ are finite constants, 
the Mittag-Leffler factor is bounded:
\begin{equation*}
    E_{\alpha}\!\left(C_1 M\,\Gamma(\alpha)\cdot T^{\alpha}\right) 
    = C_2 < \infty.
\end{equation*}
Therefore~\eqref{Tn_ML} gives:
\begin{equation}\label{Tn_a}
    |T_n(t)| \leq C_2\, a(t).
\end{equation}



By Lemma~\ref{lemm3}, the Fourier coefficients satisfy:
\begin{equation*}
    |\varphi_n| \leq \frac{C_{10}}{n^4}, 
    \qquad 
    |f_n(t)| \leq \frac{C_{11}}{n^4}, 
    \qquad n=1,2,3,\ldots
\end{equation*}
Substituting into the definition of $a(t)$:

\begin{equation*}
    a(t)\leq C_1\cdot\frac{C_{10}}{n^4}
 + C_1\cdot\frac{C_{11}}{n^4}\int_0^t(t-\tau)^{\alpha-1}\,d\tau\leq \frac{C_3}{n^{4}}, 
\end{equation*}

where $C_3 = C_1\!\left(C_{10}+\dfrac{C_{11}T^{\alpha}}{\alpha}\right)<\infty$.

Substituting the bound on $a(t)$ into~\eqref{Tn_a}:
\begin{equation*}
    |T_n(t)| \leq C_2\cdot\frac{C_3}{n^4} 
              =\frac{\widehat{C}}{n^4}, 
    \qquad t\in[0,T],\ n=1,2,3,\ldots
\end{equation*}
where $\widehat{C} = C_2 C_3$ depends only on 
$\alpha$, $M$, $T$, $l$, $\nu$ and is independent of $n$ and $t$. 
This completes the proof.


\end{proof}
We will seek a solution of the problem $\{\eqref{1.5},\eqref{2},\eqref{3}\}$in the following form:

\begin{equation}\label{4.2}
    u(t, x)=\sum_{n=1}^{+\infty}\left[\varphi_n E_{\alpha}\!\left(-\lambda_n^2 t^{\alpha}\right)
+
\int_0^t
(t-\tau)^{\alpha-1}
E_{\alpha,\alpha}\!\left(-\lambda_n^2 (t-\tau)^{\alpha}\right)
\big[q(\tau)T_n(\tau)+f_n(\tau)\big]
\, d\tau\right] X_n(x).
\end{equation}

\begin{thm}\label{10.1}
[Existence and Uniqueness]
Let the assumptions of the lemmas \ref{lemm2}, \ref{lemm3} and \ref{T(t)} be satisfied and $q(t)\in C[0,T]$. 
Then the initial-boundary value problem $\{\eqref{1.5}, \ref{2}, \eqref{3}\}$ 
admits a unique solution $u(x,t)$ such thhat: 
\begin{equation*}
    u(x, t), \quad x^\nu u_x(x, t) \in C(\bar{\Omega}), \quad x^{-\nu}\frac{\partial}{\partial x}\left[x^\nu u_x\right],\quad ^cD_{0t}^\alpha u \in C(\Omega).
\end{equation*}
\end{thm}
\begin{proof}
\textbf{ Existence:} 
To prove Theorem~\ref{10.1}, it suffices to show that the series \eqref{4.2} and the series
corresponding to the functions $x^{-\nu}\dfrac{\partial}{\partial x}\left[x^{\nu}u_x\right]$ converge uniformly in $\Omega$, and the series corresponding to the function $_{c}D_{0t}^{\alpha}u(x,t)$ converge uniformly
in any compact set $D \subset\Omega$.
First, we consider the series \eqref{4.2}. 

\begin{equation*}
    |u(x,t)|\leq \sum_{n=1}^{+\infty}|\varphi_n|\left|E_{\alpha}(-\lambda_n^{2}t^{\alpha})\right|\left|X_n(x)\right|+\sum_{n=1}^{+\infty}\int_0^t(t-\tau)^{\alpha-1}\left|E_{\alpha,\alpha}(-\lambda_n^2(t-\tau)^{\alpha})\right||X_n(x)||f_n(\tau)|d\tau
\end{equation*}
\begin{equation*}
    +\sum_{n=1}^{+\infty}\int_0^tq(\tau)(t-\tau)^{\alpha-1}\left|E_{\alpha,\alpha}\left(-\lambda_n^2(t-\tau)^{\alpha}\right)\right|\left|T_n(\tau)\right||X_n(x)|d\tau.
\end{equation*}

Using  lemmas~\ref{Mittag},\ref{lemm2} and $q(t)\in C[0,T]$. Then, we have 
\begin{equation*}
    |u(x,t)|\leq C_{5}\sum_{n=1}^{+\infty}|\varphi_n|+C_{6}\sum_{n=1}^{+\infty}\int_0^t|f_n(\tau)|(t-\tau)^{\alpha-1}d\tau+
    C_{7}\sum_{n=1}^{+\infty}\int_0^t(t-\tau)^{\alpha-1}\left|T_n(\tau)\right|d\tau
\end{equation*}

\begin{equation*}
    \leq C_{5}\sum_{n=1}^{+\infty}|\varphi_n|+C_{6}\sum_{n=1}^{+\infty}\int_0^T(t-\tau)^{\alpha-1}|f_n(\tau)|d\tau+C_{7}\sum_{n=1}^{+\infty}\int_0^T(t-\tau)^{\alpha-1}\left|T_n(\tau)\right|d\tau.
\end{equation*}
Taking into account lemma~\ref{lemm3} and lemma~\ref{T(t)}, we get the following form:
\begin{equation*}
    |u(x,t)|\leq C_{8}\sum_{n=1}^{+\infty}n^{-4}+C_{9}\sum_{n=1}^{+\infty} n^{-4}\int_0^T(t-\tau)^{\alpha-1}d\tau +C_{10}\sum_{n=1}^{+\infty}n^{-4}\int_0^T(t-\tau)^{\alpha-1}d\tau
\end{equation*}



\begin{equation*}
    \leq C_{8}\sum_{n=1}^{+\infty}n^{-4}+C_{11}\sum_{n=1}^{+\infty}n^{-4}< +\infty.
\end{equation*}

Next, we consider the series of $x^{-\nu}\dfrac{\partial}{\partial x}\left[x^{\nu}u_x\right]$,

\begin{equation*}
    \left|x^{-\nu}\dfrac{\partial}{\partial x}\left[x^{\nu}u_x\right]\right|\leq \left|\sum_{n=1}^{+\infty}x^{-\nu}\dfrac{\partial}{\partial x}\left[x^{\nu}X_n^{\prime}(x)\right]T_n(t)\right|.
\end{equation*}

Using the equation~\eqref{5} and \eqref{4.1}, we have 

\begin{equation*}
    \left|x^{-\nu}\dfrac{\partial}{\partial x}\left[x^{\nu}u_x\right]\right|\leq \left|\sum_{n=1}^{+\infty}\lambda_n^2X_n(x)T_n(t)\right|
\end{equation*}

\begin{equation*}
    \leq \sum_{n=1}^{+\infty}\left|\varphi_n E_{\alpha}\!\left(-\lambda_n^2 t^{\alpha}\right)
+
\int_0^t
(t-\tau)^{\alpha-1}
E_{\alpha,\alpha}\!\left(-\lambda_n^2 (t-\tau)^{\alpha}\right)
\big[q(\tau)T_n(\tau)+f_n(\tau)\big]
\, d\tau\right| \lambda_n^2|X_n(x)|
\end{equation*}

\begin{equation*}
    \leq\sum_{n=1}^{+\infty}\lambda_n^2\left|\varphi_n\right|\left|E_{\alpha}\left(-\lambda_n^{2}t^\alpha\right)\right||X_n(x)|+\sum_{n=1}^{+\infty}\lambda_n^2\int_0^T(t-\tau)^{\alpha-1}\left|E_{\alpha,\alpha}(-\lambda_n^2(t-\tau)^{\alpha})\right||X_n(x)||f_n(\tau)|d\tau
\end{equation*}

\begin{equation*}
    +\sum_{n=1}^{+\infty}\lambda_n^2\int_0^Tq(\tau)(t-\tau)^{\alpha-1}\left|E_{\alpha,\alpha}\left(-\lambda_n^2(t-\tau)^{\alpha}\right)\right|\left|T_n(\tau)\right||X_n(x)|d\tau.
\end{equation*}

Taking into account Lemma~\ref{Mittag}, Lemma~\ref{lemm2}, and the fact that $q(t)\in C[0,T]$, the last inequality can be written as follows.

\begin{equation*}
    \left|x^{-\nu}\dfrac{\partial}{\partial x}\left[x^{\nu}u_x\right]\right|\leq C_{12}\sum_{n=1}^{+\infty}\lambda_n^2\left|\varphi_n\right|+C_{13}\sum_{n=1}^{+\infty}\lambda_n^2\int_0^t(t-\tau)^{\alpha-1}|f_n(\tau)|d\tau+C_{14}\sum_{n=1}^{+\infty}\lambda_n^2\int_0^t(t-\tau)^{\alpha-1}\left|T_n(\tau)\right|d\tau.
\end{equation*}

Next, taking into account Lemma~\ref{lemm3}, Lemma~\ref{T(t)}, and the estimate $\lambda_n^2 \leq C n^2$, we obtain the following.


\begin{equation*}
\left|x^{-\nu}\dfrac{\partial}{\partial x}\left[x^{\nu}u_x\right]\right|\leq C_{15}\sum_{n=1}^{+\infty}n^{-2}+C_{16}\sum_{n=1}^{+\infty}n^{-2}+C_{17}\sum_{n=1}^{+\infty}n^{-2}<+\infty.
\end{equation*}

Thus, the series defining $u(x,t)$ converges absolutely and compact set $D$. Consequently, the function $u(x,t)$ satisfies the equation and the corresponding initial and boundary conditions. Hence, the problem admits a solution. 



\textbf{Uniqueness:} Suppose $u_1$ and $u_2$ are two solutions satisfying $\{\eqref{1.5},\eqref{2},\eqref{3}\}$ . Define $W(x, t)=u_1(x, t)-u_2(x, t)$. Then $W$ satisfies:
\begin{equation}\label{4.3}
    _cD_{0t}^{\alpha}W =x^{-\nu}\frac{\partial}{\partial x}\left[x^\nu W_x\right]+q(t)W(x,t), \quad(x, t) \in \Omega,
\end{equation}
\begin{equation}
    W(x, 0)  =0, \quad x \in [0,l],
\end{equation}
\begin{equation}\label{4.5}
    W(0, t)  =W(l, t)=0, \quad t \in[0, T] .
\end{equation}
Next, we will multiply equation \eqref{4.3} by $x^{\nu} W(x,t)$ and integrate over the interval $[0,l]$. As a result, we obtain the following relation:

\begin{equation*}
    \int_0^l x^{\nu}W(x,t)_cD_{0t}W(x,t)dx=\int_0^l x^{-\nu}\frac{\partial}{\partial x}\left[x^\nu W_x\right]x^{\nu}W(x,t)dx+q(t)\int_0^lx^{\nu}|W|^2dx.
\end{equation*}
Applying the integration by parts formula to the right-hand side of the last equality and using condition \eqref{4.5}, we get 
\begin{equation*}
    \int_0^l x^{\nu}W(x,t)_cD_{0t}W(x,t)dx=-\int_0^lx^{\nu}\left|W_x\right|^2dx+q(t)\int_0^lx^{\nu}|W|^2.
\end{equation*}
 Next, applying Property~1 to the last equality, we obtain
\begin{equation*}
    \frac{1}{2} {_c}D_{0t}^{\alpha}\int_0^lx^{\nu}|W(x,t)|^2dx\le-\int_0^lx^{\nu}\left|W_{x}(x,t)\right|^2dx+q(t)\int_0^lx^{\nu}|W|^2.
\end{equation*}




Using the Lemma~\ref{poincare}: 
\begin{equation*}
    \int_0^l x^{\nu}|W(x,t)|^2dx\leq C_p\int_0^l x^{\nu}|W_x(x,t)|^2dx.
\end{equation*}
Substitute into the inequality:
\begin{equation*}
    \frac{1}{2} {_c}D_{0t}^{\alpha}\int_0^lx^{\nu}|W(x,t)|^2dx\le-\frac{1}{C _p}\int_0^lx^{\nu}\left|W(x,t)\right|^2dx+q(t)\int_0^lx^{\nu}|W(x,t)|^2,
\end{equation*}

and the assumption that 
$q \in C[0,T]$ satisfies $q(t)\le M$ for all $t\in[0,T]$, we obtain
Substitute into the inequality:

\begin{equation*}
    \frac{1}{2} {_c}D_{0t}^{\alpha}\int_0^lx^{\nu}|W(x,t)|^2dx\leq-\frac{1}{C _p} \int_0^l x^\nu|W|^2 d x +M\int_0^lx^{\nu}|W(x,t)|^2.
\end{equation*}


Since $\int_0^lx^{\nu}|W|^2 d x=2 E(t)$, we get:

\begin{equation*}
    _cD_{0t}^{\alpha}E(t)\leq \left(M-\frac{1}{C _p}\right)E(t).
\end{equation*}


Let $A=M-\dfrac{1}{C _p}$. Then

\begin{equation*}
    _cD_{0t}^{\alpha}E(t) \leq A E(t).
\end{equation*}
Since $E(0)=0$ and $E(t) \geq 0$, Gronwall's inequality implies:

$$
E(t) \leq 0, \text { for all } t \in[0, T] .
$$


Therefore, $E(t)=0$ for all t , so $w(x, t)=0$ almost everywhere. Hence,

$$
u_1(x, t)=u_2(x, t) \text {, }
$$

which completes the proof.

\end{proof}




\section{ the Inverse problem 2 }
\subsection{Formulation of the problem}
In this section, within the domain $\Omega$, we examine the existence and uniqueness of solutions to the inverse source problem associated with Eq.~\eqref{1.5}.

\textbf{Inverse coefficient problem 1.}
Prove the existence and uniqueness of a pair of functions $\{u(x,t), q(t)\}$ satisfying the following conditions:
\begin{enumerate}
    \item $u(x,t),\quad x^\nu u_x(x,t) \in C(\bar{\Omega}), \quad
x^{-\nu}\frac{\partial}{\partial x}\bigl[x^\nu u_x\bigr],\quad
_
{c}\mathcal{D}_{0t}^\alpha u \in C(\Omega) ,\quad q \in C([0,T]).$
    \item They satisfy equation \eqref{1.5};
    \item $u (x, t)$ satisfies the initial condition \eqref{2}, the boundary condition \eqref{3}, and the following over-determination condition
    \begin{equation}\label{g(t)}
    \int_0^l u(x, t)x^{\nu} w(x) d x=g(t), \quad t \in[0, T].
\end{equation}
where $g(t)$, $w(x)$ are given functions.
\end{enumerate}


\subsection{ Existence and uniqueness of the solution of the inverse problem 2.}

\begin{thm}\label{T4.3}(Existence of an Inverse Problem).  Let the following conditions hold:

- $\varphi(x)$ and $f(x,t)$ satisfy the assumptions of Lemma \ref{lemm3};

- $w(x) \in C^4[0,l]$ with $w(0) = w(l) = 0$;

- $g \in C^1([0,T])$ with $g(t) \neq 0$ for all $t \in [0,T]$.

Then the inverse problem \{\eqref{1.5},\eqref{2},\eqref{3},\eqref{g(t)}\} admits at least one solution pair
$\{q(t),\, u(x,t)\}$ such that

\begin{equation*}
    q \in C([0,T]),\quad u(x,t),\quad x^\nu u_x(x,t) \in C(\bar{\Omega}), \quad
x^{-\nu}\frac{\partial}{\partial x}\bigl[x^\nu u_x\bigr],\quad
_
{c}\mathcal{D}_{0t}^\alpha u \in C(\Omega).
\end{equation*}                                               
\end{thm} 

\begin{proof}
 
 Given an arbitrary $\hat{q} \in C([0,T])$, it follows from Theorem~\ref{10.1} that
the direct problem $\{\eqref{1.5},\eqref{2},\eqref{3}\}$ possesses a unique solution
$u_{\hat{q}}(x,t)$ with the regularity
\begin{equation*}
    u_{\hat{q}},\quad x^\nu (u_{\hat{q}})_x \in C(\bar{\Omega}), \qquad
    x^{-\nu}\frac{\partial}{\partial x}\bigl[x^\nu (u_{\hat{q}})_x\bigr],\quad
    {}_{c}\mathcal{D}_{0t}^\alpha u_{\hat{q}} \in C(\Omega).
\end{equation*}
Our argument is based on Schauder's fixed-point theorem (Theorem~\ref{Schauder}).
To this end, we introduce the operator $\mathcal{A}: C([0,T])\to C([0,T])$ as follows:
\begin{equation}
\label{operator_A}
    (\mathcal{A}\hat{q})(t) :=
    \frac{{}_{c}\mathcal{D}_{0t}^\alpha g(t)
    +\displaystyle\int_0^l x^\nu (u_{\hat{q}})_x\,w_x\,dx
    -\displaystyle\int_0^l x^\nu f(x,t)\,w(x)\,dx}{g(t)}.
\end{equation}
Here the boundary contributions $\bigl[x^\nu u_{\hat{q}} w_x\bigr]_0^l$ and
$-\bigl[x^\nu (u_{\hat{q}})_x w\bigr]_0^l$ are zero, which follows from the conditions
$w(0)=w(l)=0$ and $u_{\hat{q}}(0,t)=u_{\hat{q}}(l,t)=0$.

\textbf{Continuity.}
Suppose $\hat{q}_n\to\hat{q}$ in $C([0,T])$.
The continuous dependence result of Theorem~\ref{10.1} yields that
both $u_{\hat{q}_n}\to u_{\hat{q}}$ and
$x^\nu(u_{\hat{q}_n})_x \to x^\nu(u_{\hat{q}})_x$
hold uniformly on $\bar{\Omega}$.
As each term in~\eqref{operator_A} converges uniformly,
we conclude that $\mathcal{A}\hat{q}_n\to\mathcal{A}\hat{q}$ in $C([0,T])$,
establishing continuity of $\mathcal{A}$.

\textbf{Compactness.}
Let $P>0$ be given and set $K=\{q\in C([0,T]):\|q\|_{C([0,T])}\leq P\}$.
According to Theorem~\ref{10.1}, there exists a constant $C_0>0$,
depending solely on $\varphi$, $f$, $l$, $T$ and in particular
independent of $\hat{q}\in K$, such that
\begin{equation*}
    \sup_{(x,t)\in\bar{\Omega}}\bigl|x^\nu(u_{\hat{q}})_x(x,t)\bigr| \leq C_0,
    \qquad \forall\, \hat{q}\in K.
\end{equation*}
Consequently, every $\hat{q}\in K$ satisfies the estimate
\begin{equation*}
    \|(\mathcal{A}\hat{q})\|_{C([0,T])}
    \leq
    \frac{1}{\displaystyle\min_{t\in[0,T]}|g(t)|}
    \Bigl(
        \|{}_{c}\mathcal{D}_{0t}^\alpha g\|_{C([0,T])}
        + C_0\|w_x\|_{C([0,l])}
        + \|f\|_{C(\bar{\Omega})}\|w\|_{C([0,l])}
    \Bigr) =: P.
\end{equation*}
Moreover, since $x^\nu(u_{\hat{q}})_x\in C(\bar{\Omega})$ uniformly over $K$,
the collection $\{\mathcal{A}\hat{q}:\hat{q}\in K\}$ is both equicontinuous
and uniformly bounded in $C([0,T])$.
Hence $\mathcal{A}(K)$ is relatively compact in $C([0,T])$
by the Arzel\`{a}--Ascoli theorem (Theorem~\ref{Ascoli}).

\textbf{Invariance.}
The set $K$ is a nonempty, closed, bounded, and convex subset of the
Banach space $C([0,T])$, and the estimate above shows that $\mathcal{A}(K)\subseteq K$.

An application of Schauder's fixed-point theorem (Theorem~\ref{Schauder})
now yields a fixed point $q\in K\subset C([0,T])$ of $\mathcal{A}$.
Setting $u:=u_q$, the pair $\{q,u\}$ is a solution of
$\{\eqref{1.5},\eqref{2},\eqref{3},\eqref{g(t)}\}$ satisfying
\begin{equation*}
    q\in C([0,T]),\qquad
    u,\;x^\nu u_x\in C(\bar{\Omega}),\qquad
    x^{-\nu}\frac{\partial}{\partial x}[x^\nu u_x],\;
    {}_{c}\mathcal{D}_{0t}^\alpha u\in C(\Omega).
\end{equation*}
\end{proof}


\begin{thm}
(Uniqueness of the Solution). Let the assumptions of Lemma~\ref{lemm3} hold, and suppose in addition that
$w \in C^1([0,l])$ with $w(0)=w(l)=0$ and $g \in C^1([0,T])$ with $g(t) \neq 0$ for all $t \in [0,T]$.
Then the solution pair $\{q(t), u(x,t)\}$ to the inverse problem
\{\eqref{1.5},\eqref{2},\eqref{3},\eqref{g(t)}\} is unique.
\end{thm} 
\begin{proof}    
 Suppose $\{q, u\}$ and $\{\tilde{q}, \tilde{u}\}$ are solution pairs. Define

$$
U(x, t):=u(x, t)-\tilde{u}(x, t), \quad R(t):=q(t)-\tilde{q}(t) .
$$


Subtracting the two PDEs gives:

\begin{equation}\label{4.6}
    _{c}\mathcal{D}_{0t}^{\alpha}U =x^{-\nu}\frac{\partial}{\partial x}\left[x^{\nu}U_x\right]+q(t)U+R(t)\tilde{u} \quad(x, t) \in \Omega, 
\end{equation}
\begin{equation*}
    U(x, 0)  =0, \quad x \in[0, l], 
\end{equation*}
\begin{equation*}
    U(0, t)  =U(l, t)=0, \quad t \in[0, T],
\end{equation*}
Applying the operator ${_c}D_{0t}^{\alpha}$ to relation \eqref{g(t)} corresponding to each pair of solutions $(u,\tilde{u})$, we get :

\begin{equation*}
    \int_0^l  w(x){_c}\mathcal{D}_{0t}^{\alpha}u d x={_c}\mathcal{D}_{0t}^{\alpha}g(t), \quad \int_0^l w(x){_c}\mathcal{D}_{0t}^{\alpha}\tilde{u}_t  d x={_c}\mathcal{D}_{0t}^{\alpha}g(t),\quad t \in[0, T]
\end{equation*}

with $g$ from \eqref{g(t)}.
Substitute the PDE \eqref{1.5} into these expressions. Using integration by parts and the boundary conditions $w(0)=w(l)=0$, we obtain 


\begin{equation*}
    {_c}\mathcal{D}_{0t}^{\alpha}g(t)=-\int_0^l x^{\nu}u_{x} w_x(x) d x+q(t)g(t)+\int_0^l x^{\nu}f(x,t) w(x) d x 
\end{equation*}

\begin{equation*}
     {_c}\mathcal{D}_{0t}^{\alpha}g(t)=-\int_0^l x^{\nu}\tilde{u}_{x} w_x(x) d x+\tilde{q}(t)g(t)+\int_0^l x^{\nu}f(x,t) w(x) d x 
\end{equation*}



Subtracting the two identities yields

\begin{equation*}
    0=-\int_0^l x^{\nu}U_{x}w_x(x)dx+R(t)g(t)
\end{equation*}

Since $g(t) \neq 0$ on $[0, T]$, we can write

\begin{equation}\label{4.7}
    R(t)=\frac{\int_0^l x^{\nu}U_{x}w_x(x)dx}{g(t)}
\end{equation}
Next, we multiply equation \eqref{4.6} by $x^{\nu}U$ and integrate it with respect to the variable $x$ over the interval $[0,l]$. As a result, we obtain 

\begin{equation*}
    \int_0^l x^{\nu}U{_c}\mathcal{D}_{0t}^{\alpha}Udx=\int_0^l U\frac{\partial}{\partial x}\left[x^{\nu}U_x\right]dx+q(t)\int_0^lx^{\nu}|U|^2dx+R(t)\int_0^l x^{\nu}U\tilde{u}dx
 \end{equation*}



Using the Property 1, we have 

\begin{equation*}
    \frac{1}{2}{_c}\mathcal{D}_{0t}^\alpha \int_0^l x^{\nu} |U|^2\leq \int_0^l U\frac{\partial}{\partial x}\left[x^{\nu}U_x\right]dx+q(t)\int_0^lx^{\nu}|U|^2dx+R(t)\int_0^l x^{\nu}U\tilde{u}dx
 \end{equation*}




and Consider the energy functional $E(t)=\frac{1}{2} \int_0^lx^{\nu}|U|^2 d x$.

\begin{equation*}
    {_c}\mathcal{D}_{0t}^{\alpha}E(t)\leq \underbrace{\int_0^l U\frac{\partial}{\partial x}\left[x^{\nu}U_x\right]dx}_{I_1}+q(t)\underbrace{\int_0^lx^{\nu}|U|^2dx}_{I_2}+R(t)\underbrace{\int_0^l x^{\nu}U\tilde{u}dx}_{I_3}
\end{equation*}


For $I_1$, integrate by parts and apply boundary conditions $U(0, t)=U(l, t)=0$ :

\begin{equation*}
    I_1=\left[x^{\nu}U U_x\right]_0^l-\int_0^lx^{\nu}\left|U_x\right|^2 d x=-\left\|U_x\right\|_{L_{\nu}^2}^2
\end{equation*}

Thus:

\begin{equation}\label{4.8}
     {_c}\mathcal{D}_{0t}^{\alpha} E(t)\leq-\left\|U_x\right\|_{L_{\nu}^2}^2+q(t)\|U\|_{L_{\nu}^2}^2+R(t) \int_0^l x^{\nu}U \tilde{u} d x.
\end{equation}



From \eqref{4.7} and the Cauchy-Schwarz inequality:

\begin{equation*}
    |R(t)| \leq \frac{1}{|g(t)|}\left|\int_0^l x^{\nu}U_x w_{x}(x) d x\right| \leq \frac{\left\|w_{x}\right\|_{L_{\nu}^2}}{|g(t)|}\|U_x\|_{L_{\nu}^2}
\end{equation*}


Define the following constants:
\begin{equation*}
 M_1 =\sup _{t \in[0, T]}\|\tilde{u}(\cdot, t)\|_{L_{\nu}^{2}}<\infty \quad\left(\text { since } \tilde{u}(x,t), \quad x^{\nu}\tilde{u}_x\in C(\overline{\Omega})\right.  \text { by Theorem \ref{10.1}) }   
\end{equation*}

\begin{equation*}
    M_2 =\left\|w_{x}\right\|_{L_{\nu}^{2
    }}<\infty \quad\left(\text { since } w(x) \in C^1([0, l])\right. 
\end{equation*}

\begin{equation*}
    m =\min _{t \in[0, T]}|g(t)|>0 \quad \text { (since } g \in C^1([0, T]) \text { and } g(t) \neq 0).
\end{equation*}

Then:

\begin{equation}\label{4.9}
    |R(t)| \leq \frac{M_2}{m}\|U_x\|_{L_{\nu}^2}
\end{equation}


Applying Cauchy-Schwarz to $I_3$ :

\begin{equation}\label{4.10}
    \left| \int_0^l x^{\nu}U \tilde{u} d x\right| \leq\|U\|_{L_{\nu}^2}\|\tilde{u}\|_{L_{\nu}^2} \leq M_1\|U\|_{L_{\nu}^2} .
\end{equation}


Substitute \eqref{4.9} and \eqref{4.10} into \eqref{4.8}:

\begin{equation*}
 {_c}\mathcal{D}_{0t}^{\alpha} E(t)  \leq -\left\|U_x\right\|_{L_{\nu}^2}^2+q(t)\|U\|_{L_{\nu}^2}^2+|R(t)| M_1\|U\|_{L_{\nu}^2}   
\end{equation*}
\begin{equation*}
    \leq -\left\|U_x\right\|_{L_{\nu}^2}^2+q(t)\|U\|_{L_{\nu}^2}^2+\left(\frac{M_2}{m}\|U_x\|_{L_{\nu}^2}\right) M_1\|U\|_{L_{\nu}^2}.
\end{equation*}
Apply the Young's inequality, 

\begin{equation*}
    \|U_x\|_{L_{\nu}^2}\|U\|_{L_{\nu}^2}\leq \varepsilon\|U_x\|_{L_{\nu}^2}^2+\frac{1}{4\varepsilon}\|U\|_{L_{\nu}^2}^2
\end{equation*}

Then. We can rewrite 
\begin{equation*}
    {_c}\mathcal{D}_{0t}^{\alpha} E(t)  \leq -\left\|U_x\right\|_{L_{\nu}^2}^2+q(t)\|U\|_{L_{\nu}^2}^2+\frac{M_1M_2}{m}\varepsilon\|U_x\|_{L_{\nu}^2}^2+\frac{M_1M_2}{4m\varepsilon}\|U\|_{L_{\nu}^2}^2
\end{equation*}
\begin{equation*}
    =-\left(1-\frac{M_1M_2}{m}\varepsilon\right)\left\|U_x\right\|_{L_{\nu}^2}^2+\left(q(t)+\frac{M_1M_2}{4m\varepsilon}\right)\|U\|_{L_{\nu}^2}^2
\end{equation*}
Since $\varepsilon>0$ is arbitrary, we choose it sufficiently small so that
\begin{equation*}
    1-\frac{M_1M_2}{m}\varepsilon>0.
\end{equation*}
For instance, taking
\begin{equation*}
    \varepsilon=\frac{m}{2M_1M_2},
\end{equation*}
we obtain
\begin{equation*}
    1-\frac{M_1M_2}{m}\varepsilon
=
\frac12>0.
\end{equation*}
Thus, the coefficient of $\|U_x\|_{L_{\nu}^2}^2$ remains negative. 
Discarding this non–positive term, we arrive at
\begin{equation*}
    {_c}\mathcal{D}_{0t}^{\alpha}E(t)
\le
\left(q(t)+\frac{(M_1M_2)^2}{2m^2}\right)\|U\|_{L_{\nu}^2}^2.
\end{equation*}

Since $\|U\|_{L_{\nu}^2}^2=2 E(t)$, we obtain:

\begin{equation*}
     {_c}\mathcal{D}_{0t}^{\alpha} E(t) \leq 2\left(q(t)+\frac{(M_1M_2)^2}{2m^2}\right) E(t)
\end{equation*}


Let $K:=\sup _{t \in[0, T]}\left|\left(q(t)+\frac{(M_1M_2)^2}{2m^2}\right)\right|<\infty$ (finite because $q \in C([0, T])$ and $M_1, M_2, m$ are finite with $m>0$. Then:

\begin{equation}\label{4.11}
    {_c}D_{0t}^{\alpha}E(t) \leq 2 K E(t) .
\end{equation}

Since $E(0)=0$, Grönwall's inequality applied to \eqref{4.11} implies:

\begin{equation*}
    E(t)=0 \quad \forall t \in[0, T] .
\end{equation*}


Hence, $U=0$ for all $(x, t) \in \bar{\Omega}$, so $u=\tilde{u}$. Substituting $U=0$ into \eqref{4.7} gives:

\begin{equation*}
    R(t) g(t)=0 \Longrightarrow R(t)=0 \quad \forall t \in[0, T],
\end{equation*}


Since $g(t) \neq 0$, thus, $q=\tilde{q}$, proving uniqueness.
\end{proof}

\newpage
































































\newpage
\begin{thebibliography}{99}

\bibitem{Carslaw} H.S.~Carslaw, J.C.~Jaeger,
\textit{Conduction of Heat in Solids},
2nd ed., Clarendon Press, Oxford, 1959.

\bibitem{bib6}  Erd´elyi~A., Magnus~W., Oberhettinger~F., Tricomi~F.G.; Higher transcendental functions, vol. 2, Based, in part, on notes left by H. Bateman,
McGraw-Hill, New York–Toronto–London 1953; Russian transl., Nauka, Moscow
1966.

\bibitem{Watson}  Watson~G.N.; A treatise on the theory of Bessel functions, 2nd ed., Cambridge Univ. Press, Cambridge; Macmillan, New York 1944; Russian transl., Part 1, Moscow, IL 1949.

\bibitem{Olver} Olver~F.W.J.; Introduction to asymptotics and special functions, Academic Press, New York 1974; Russian transl., Nauka, Moscow 1978. 

\bibitem{Sabitov} Sabitov~K.B.,  Suleimanova~A. Kh.; “The Dirichlet problem for a mixed-type equation of the second kind in a rectangular domain”, Izv. Vyssh. Ucheb. Zaved. Mat., 2007, no. 4, 45–53; English transl., Russian Math. (Iz. VUZ ) 51:4 (2007),42–50.
\bibitem{Ye2007}
H. Ye, J. Gao, Y. Ding, A generalized Gronwall inequality 
and its application to a fractional differential equation, 
J. Math. Anal. Appl. 328 (2007), 1075–1081.
\bibitem{Vagapova} Sabitov~K. B., Vagapova~ E. V.; “Dirichlet problem for an equation of mixed type with two degeneration lines in a rectangular domain”, Differentsial’nye Uravneniya 49:1 (2013), 68–78; English transl., Differ. Equ. 49:1 (2013), 68–78.


\bibitem{lemma3.2}K. B. Sabitov, R. M. Safina, The first boundary-value problem for an equation of mixed type with a singular coefficient, Izvestiya: Mathematics, 2018, Volume 82, Issue 2, 318–350

\bibitem{Pedas}A.Pedas, G.Vainikko. Integral equations with diagonal and boundary singularities of the kernel. 2006.
Z. Anal. Anwend, pp. 487–516.

\bibitem{Zeidler}Zeidler E., Nonlinear Functional Analysis and its Applications I: Fixed-Point Theorems,
Springer-Verlag, New York, 1986.

\bibitem{Podlubny} Podlubny, I.: Fractional Differential Equations: An Introduction to Fractional Derivatives, Fractional
Differential Equations, to Methods of Their Solution and Some of Their Applications. Mathematics in
Science and Engineering, vol. 198, Academic Press, San Diego (1999)

\bibitem{Evans}Evans L.C., Partial Differential Equations, 2nd ed., Graduate Studies in Mathematics, Vol. 19, American Mathematical Society, Providence, RI, 2010. 19, American Mathematical Society, Providence, RI, 2010.

\bibitem{A1} M.\,S.~Hussein, D.~Lesnic, V.\,L.~Kamynin, A.\,B.~Kostin,
   Direct and inverse source problems for degenerate parabolic equations,
   \textit{J. Inverse Ill-Posed Problems}, \textbf{28} (2020), 425--448.

 \bibitem{A2} P.~Cannarsa, J.~Tort, M.~Yamamoto,
   Determination of source terms in a degenerate parabolic equation,
   \textit{Inverse Problems}, \textbf{26} (2010), 105003.

 \bibitem{A3} V.\,L.~Kamynin,
   On inverse problems for strongly degenerate parabolic equations
   under the integral observation condition,
   \textit{Comput. Math. Math. Phys.}, \textbf{58} (2018), 2002--2017.

 \bibitem{A4} N.~Huzyk, P.~Pukach, M.~Vovk,
   Coefficient inverse problem for the strongly degenerate parabolic equation,
   \textit{Carpathian Math. Publ.}, \textbf{15} (2023), 52--65.

 \bibitem{A5} Y.~Luchko,
   Some uniqueness and existence results for the initial-boundary-value problems
   for the generalized time-fractional diffusion equation,
   \textit{Comput. Math. Appl.}, \textbf{59} (2010), 1766--1772.

 \bibitem{A6} P.~Agarwal, E.~Karimov, M.~Mamchuev, M.~Ruzhansky,
   On boundary-value problems for a partial differential equation
   with Caputo and Bessel operators,
   \textit{Springer}, 2017, 127--140.

 \bibitem{A7} D.\,K.~Durdiev,
   Inverse coefficient problem for the time-fractional diffusion equation
   with Hilfer operator,
   \textit{Math. Methods Appl. Sci.}, \textbf{46} (2023), 11099--11108.

 \bibitem{A8} K.~Sakamoto, M.~Yamamoto,
   Inverse source problem with a final overdetermination for a fractional
   diffusion equation,
   \textit{Math. Control Relat. Fields}, \textbf{1} (2011), 509--518.
\ Anal.\ Appl.\ \textbf{328} (2007), 1075--1081.
\bibitem{Akramova}Akramova D.I., Inverse coefficient problem for a fractional-diffusion equation with a Bessel operator, Izvestiya vuzov. Matematika, 2023, no. 9, 45–57.
\end{thebibliography}



\label{lastpage}

%\endinput

%
\end{document}
